% Options for packages loaded elsewhere
\PassOptionsToPackage{unicode}{hyperref}
\PassOptionsToPackage{hyphens}{url}
\PassOptionsToPackage{dvipsnames,svgnames,x11names}{xcolor}
%
\documentclass[
]{article}
\usepackage{amsmath,amssymb}
\usepackage{iftex}
\ifPDFTeX
  \usepackage[T1]{fontenc}
  \usepackage[utf8]{inputenc}
  \usepackage{textcomp} % provide euro and other symbols
\else % if luatex or xetex
  \usepackage{unicode-math} % this also loads fontspec
  \defaultfontfeatures{Scale=MatchLowercase}
  \defaultfontfeatures[\rmfamily]{Ligatures=TeX,Scale=1}
\fi
\usepackage{lmodern}
\ifPDFTeX\else
  % xetex/luatex font selection
\fi
% Use upquote if available, for straight quotes in verbatim environments
\IfFileExists{upquote.sty}{\usepackage{upquote}}{}
\IfFileExists{microtype.sty}{% use microtype if available
  \usepackage[]{microtype}
  \UseMicrotypeSet[protrusion]{basicmath} % disable protrusion for tt fonts
}{}
\makeatletter
\@ifundefined{KOMAClassName}{% if non-KOMA class
  \IfFileExists{parskip.sty}{%
    \usepackage{parskip}
  }{% else
    \setlength{\parindent}{0pt}
    \setlength{\parskip}{6pt plus 2pt minus 1pt}}
}{% if KOMA class
  \KOMAoptions{parskip=half}}
\makeatother
\usepackage{xcolor}
\usepackage[margin = 0.5in]{geometry}
\usepackage{graphicx}
\makeatletter
\def\maxwidth{\ifdim\Gin@nat@width>\linewidth\linewidth\else\Gin@nat@width\fi}
\def\maxheight{\ifdim\Gin@nat@height>\textheight\textheight\else\Gin@nat@height\fi}
\makeatother
% Scale images if necessary, so that they will not overflow the page
% margins by default, and it is still possible to overwrite the defaults
% using explicit options in \includegraphics[width, height, ...]{}
\setkeys{Gin}{width=\maxwidth,height=\maxheight,keepaspectratio}
% Set default figure placement to htbp
\makeatletter
\def\fps@figure{htbp}
\makeatother
\setlength{\emergencystretch}{3em} % prevent overfull lines
\providecommand{\tightlist}{%
  \setlength{\itemsep}{0pt}\setlength{\parskip}{0pt}}
\setcounter{secnumdepth}{-\maxdimen} % remove section numbering
\usepackage{color}
\usepackage{framed}
\setlength{\fboxsep}{.8em}

\usepackage{tcolorbox}

\tcbset{
  mystyle/.style={
    before skip=-10pt, % Adjust this value
    after skip=-10pt   % Adjust this value
  }
}
\newtcolorbox{mynoteblock}[1][]{mystyle, #1}

\newtcolorbox{blackbox}{
  colback=black,
  colframe=orange,
  coltext=white,
  boxsep=5pt,
  arc=4pt}
  
\newtcolorbox{ltbluebox}{
  colback=blue!5!white,
  colframe=blue!75!black,
  boxsep=5pt,
  arc=4pt}
  
\newtcolorbox{ltgreenbox}{
  colback=green!5!white,
  colframe=blue!75!black,
  boxsep=5pt,
  arc=4pt}
  
\newtcolorbox{ltpurplebox}{
  colback=purple!5!white,
  colframe=blue!75!black,
  boxsep=5pt,
  arc=4pt}
  


\newenvironment{cols}[1][]{}{}

\newenvironment{col}[1]{\begin{minipage}[t]{#1}\ignorespaces\vspace{0pt}}{%
\end{minipage}
\ifhmode\unskip\fi
\aftergroup\useignorespacesandallpars}

\def\useignorespacesandallpars#1\ignorespaces\fi{%
#1\fi\ignorespacesandallpars}

\makeatletter
\def\ignorespacesandallpars{%
  \@ifnextchar\par
    {\expandafter\ignorespacesandallpars\@gobble}%
    {}%
}
\makeatother

\usepackage{awesomebox}
\setlength{\aweboxvskip}{1mm}


\pagenumbering{gobble}

\usepackage{enumitem}
\setlist{itemsep=1pt, parsep=0pt} %# Adjust vertical spacing
\setlist[itemize,1]{leftmargin=*, labelsep=0.3em} %# Adjust horizontal spacing
\setlist[enumerate,1]{leftmargin=*, labelsep=0.3em} %# Adjust horizontal spacing

\ifLuaTeX
  \usepackage{selnolig}  % disable illegal ligatures
\fi
\usepackage{bookmark}
\IfFileExists{xurl.sty}{\usepackage{xurl}}{} % add URL line breaks if available
\urlstyle{same}
\hypersetup{
  colorlinks=true,
  linkcolor={Maroon},
  filecolor={Maroon},
  citecolor={Blue},
  urlcolor={blue},
  pdfcreator={LaTeX via pandoc}}

\author{}
\date{\vspace{-2.5em}}

\begin{document}

\vspace{-1cm}

\section{Measures Setting: Potential Indicators and
Pathways}\label{measures-setting-potential-indicators-and-pathways}

\begin{noteblock}
Measure setting is the process of establishing regulatory controls and
conservation tools that govern how, when, where, and how much fish can
be harvested to prevent overfishing and ensure the long-term
sustainability of stocks. Measures could include regulations such as
effort controls, spatial or temporal closures, or size and age limits to
landings. \textbf{Note: }measure setting also includes regulating catch
limits, however due to the importance of catch management decisions,
catch and quota setting are described in their own 1-pager.

\end{noteblock}

\begin{itemize}
\tightlist
\item
  \textbf{Key question: What additional indicators could inform measures
  setting? Could ecosystem indicators or recreational demand indicators
  in addition to current information improve the process?}
\item
  \textbf{Scope: Should we develop a general process for using
  indicators in the measures setting process, or focus on a specific
  application for a single species? Which species?}
\end{itemize}

\begin{cols}

\begin{col}{0.49\textwidth}

\begin{ltbluebox}

\subsection{What information is used
now?}\label{what-information-is-used-now}

Current measures in place are outlined by NOAA's
\href{https://www.fisheries.noaa.gov/species-directory}{Specification
Summaries} and their
\href{https://www.fisheries.noaa.gov/rules-and-regulations/fisheries}{Federal
Regulations}.

\subsection{Information sources for
indicators}\label{information-sources-for-indicators}

Big picture performance metrics for commercial and recreational
fisheries are available in:

\begin{itemize}
\tightlist
\item
  \href{https://www.fisheries.noaa.gov/new-england-mid-atlantic/ecosystems/state-ecosystem-reports-northeast-us-shelf}{State
  of the Ecosystem Reports}
\item
  \href{https://static1.squarespace.com/static/511cdc7fe4b00307a2628ac6/t/672bbef75e7b1166ec15dba9/1730920200863/2_MAB_RiskAssess_2024.pdf}{EAFM
  Risk Assessments}
\item
  \href{https://www.fisheries.noaa.gov/new-england-mid-atlantic/ecosystems/ecosystem-and-socioeconomic-profile-development-and-reports\#longfin-inshore-squid\%E2\%80\%94in-process}{Ecosystem
  and Socioeconomic Profiles}
\end{itemize}

Also see \href{https://www.rhapcast.org/}{river herring} bycatch
mitigation IRA project

\end{ltbluebox}

\end{col}

\begin{col}{0.02\textwidth}
~

\end{col}

\begin{col}{0.49\textwidth}

\begin{ltgreenbox}

\subsection{Information gaps}\label{information-gaps}

To understand multispecies, multifleet implications of measures set for
a given species, we would need easy access to:

\begin{itemize}
\item
  What fleets the species is caught in.

  \begin{itemize}
  \tightlist
  \item
    \href{https://www.fisheries.noaa.gov/new-england-mid-atlantic/commercial-fishing/quota-monitoring-greater-atlantic-region}{CAMS}
    data for commercial fleets
  \item
    \href{https://www.fisheries.noaa.gov/recreational-fishing-data/introduction-marine-recreational-information-program-data}{MRIP}
    for recreational fleets
  \end{itemize}
\item
  Other targets and catch from these fleets
\item
  Who manages the other species and fleets if not MAFMC
\item
  Measures in place for all the other species (commercial and
  recreational) caught with the species of interest
\end{itemize}

\textbf{Further useful information could include:}

\begin{itemize}
\item
  Aggregate performance metrics for the fisheries to evaluate
  implications of changes (commercial and recreational):

  \begin{itemize}
  \tightlist
  \item
    Landings
  \item
    Revenues
  \item
    Trips
  \item
    Expenditures
  \item
    Resilience
  \item
    Other system level fishery elements from risk assessment
  \end{itemize}
\item
  Stock projections of age and size compositions
\item
  Projections of commercial fleet and angler response to measures

  \begin{itemize}
  \tightlist
  \item
    Is there a threshold for switching to a different species?
  \end{itemize}
\item
  Projections of climate and habitat conditions for relevant species

  \begin{itemize}
  \tightlist
  \item
    See \href{https://psl.noaa.gov/cefi_portal/}{Changing Ecosystems and
    Fisheries Initiative Portal}
  \end{itemize}
\end{itemize}

\end{ltgreenbox}

\end{col}

\end{cols}

\end{document}
